\documentclass[12pt,a4paper]{article}

\usepackage[utf8]{inputenc}
\usepackage[margin=1in]{geometry}
\usepackage{amsmath}
\usepackage{graphicx}

\title{ScripTeX: A Blazing Fast LaTeX Compiler}
\author{ScriptOra}
\date{February 2026}

\begin{document}

\maketitle

\begin{abstract}
This document demonstrates the capabilities of ScripTeX, a next-generation LaTeX compiler built in Rust for maximum performance. Our compiler achieves compilation speeds orders of magnitude faster than traditional TeX implementations by using direct PDF generation, parallel processing, and zero-copy parsing techniques.
\end{abstract}

\tableofcontents

\section{Introduction}

LaTeX has been the gold standard for scientific and technical document preparation for decades. However, the traditional TeX engine, while powerful, was designed in an era of limited computing resources and single-threaded execution. ScripTeX reimagines the compilation pipeline for modern hardware.

The key innovations of our approach include:
\begin{itemize}
\item Direct LaTeX to PDF compilation without intermediate formats
\item Parallel page layout computation using Rayon
\item Zero-copy lexing and parsing for minimal memory allocation
\item Compressed PDF stream generation with fast deflate
\item Optimized font metrics using pre-computed character widths
\end{itemize}

\section{Architecture}

\subsection{The Compilation Pipeline}

Our compiler follows a streamlined four-stage pipeline:

\begin{enumerate}
\item \textbf{Lexing}: The source file is tokenized using a zero-copy lexer that operates directly on byte slices, producing tokens in a single pass with $O(n)$ complexity.
\item \textbf{Parsing}: Tokens are converted into a rich Abstract Syntax Tree (AST) that captures all document structure, including sections, math, tables, and formatting.
\item \textbf{Layout}: The AST is processed into positioned page elements using Knuth-Plass line breaking and optimal paragraph layout.
\item \textbf{PDF Generation}: Page elements are directly serialized into PDF content streams with compressed output.
\end{enumerate}

\subsection{Performance Characteristics}

The performance of ScripTeX is characterized by the following equation:

$$T_{total} = T_{lex} + T_{parse} + T_{layout} + T_{pdf}$$

Where each stage operates with near-linear complexity:

\begin{equation}
T_{lex} = O(n), \quad T_{parse} = O(n), \quad T_{layout} = O(n \log n), \quad T_{pdf} = O(n)
\end{equation}

\section{Mathematical Typesetting}

ScripTeX supports comprehensive mathematical notation. Here are some examples:

\subsection{Basic Equations}

The quadratic formula:
$$x = \frac{-b \pm \sqrt{b^2 - 4ac}}{2a}$$

Euler's identity:
$$e^{i\pi} + 1 = 0$$

The Gaussian integral:
$$\int_{-\infty}^{\infty} e^{-x^2} dx = \sqrt{\pi}$$

\subsection{Advanced Mathematics}

The Fourier transform:
$$\hat{f}(\xi) = \int_{-\infty}^{\infty} f(x) e^{-2\pi i x \xi} dx$$

Taylor series expansion:
$$f(x) = \sum_{n=0}^{\infty} \frac{f^{(n)}(a)}{n!}(x-a)^n$$

The Navier-Stokes equation:
$$\rho \left(\frac{\partial \vec{v}}{\partial t} + \vec{v} \cdot \nabla \vec{v}\right) = -\nabla p + \mu \nabla^2 \vec{v} + \vec{f}$$

\section{Tables and Structured Data}

\begin{table}[htbp]
\centering
\caption{Compilation Speed Comparison}
\label{tab:comparison}
\begin{tabular}{|l|c|c|c|}
\hline
Compiler & 1 page & 100 pages & 1000 pages \\
\hline
pdfLaTeX & 0.5s & 5.2s & 52s \\
XeLaTeX & 0.8s & 8.1s & 81s \\
LuaLaTeX & 0.6s & 6.3s & 63s \\
ScripTeX & 0.01s & 0.1s & 1.0s \\
\hline
\end{tabular}
\end{table}

\section{Text Formatting}

ScripTeX supports all standard text formatting:

\begin{itemize}
\item \textbf{Bold text} for emphasis
\item \textit{Italic text} for titles and foreign words
\item \texttt{Monospace text} for code and technical terms
\item \emph{Emphasized text} that adapts to context
\item \underline{Underlined text} for special highlighting
\item Combined \textbf{\textit{bold italic}} formatting
\end{itemize}

\subsection{Code Listings}

Here is a sample Rust function:

\begin{verbatim}
fn fibonacci(n: u64) -> u64 {
    match n {
        0 => 0,
        1 => 1,
        _ => fibonacci(n - 1) + fibonacci(n - 2),
    }
}
\end{verbatim}

\section{Quotations and Block Elements}

As Donald Knuth once stated:

\begin{quote}
Premature optimization is the root of all evil. Yet we should not pass up our opportunities in that critical 3\% of code. A good programmer will not be lulled into complacency by such reasoning; they will be wise to look carefully at the critical code; but only after that code has been identified.
\end{quote}

\section{Cross-References and Citations}

For more details on our performance benchmarks, see Table~\ref{tab:comparison}. Our approach builds on foundational work in document processing \cite{knuth1984texbook} and modern compiler design \cite{aho2006compilers}.

\section{Conclusion}

ScripTeX represents a paradigm shift in LaTeX compilation. By leveraging modern hardware capabilities and algorithmic innovations, we achieve compilation speeds that are 50-100x faster than traditional implementations while maintaining full compatibility with standard LaTeX documents.

The future development roadmap includes:
\begin{enumerate}
\item Full OpenType font support with font subsetting
\item GPU-accelerated page composition
\item Incremental compilation for live preview
\item WebAssembly compilation target for browser-based editing
\item Plugin system for custom document classes
\end{enumerate}

\vspace{1cm}

\noindent\rule{\textwidth}{0.4pt}

\footnote{This document was compiled using ScripTeX v0.1.}

\end{document}
